\section{引言}

函数逼近论研究一个基本问题：如何用简单函数（多项式、三角多项式等）去近似复杂的连续函数？这一问题贯穿数学分析的诸多分支，Weierstrass 的两个逼近定理从原则上肯定了连续函数可用多项式或三角多项式一致逼近，但并未给出具体的构造方法。

逼近论的发展由此分化为两个相互关联的研究方向。

其一是\textbf{最佳逼近理论}。给定逼近工具的“次数”约束，是否存在一个最优的逼近元，使得逼近误差达到理论上的下确界？这个最优逼近元具有怎样的特征？切比雪夫交错定理对这一问题给出了精确的刻画——误差函数在足够多的点上达到最大值且符号交替。最佳逼近理论确立了一个不可逾越的“最优基准”，但除了少数特例，直接求出最佳逼近多项式往往极其困难。

其二是\textbf{构造性逼近理论}。能否设计结构简单、便于计算的线性算子序列，使之能一致逼近任意连续函数？Bernstein 多项式、Fejér 算子等都是这一思路的经典产物。这些算子虽然未必是最优的，但胜在形式明确、可直接计算。

由此自然产生两个核心问题：
\begin{enumerate}
    \item 对于给定的线性算子序列，如何简单有效地判定它是否对一切连续函数一致收敛？
    \item 构造出来的线性算子，其逼近效果能有多好——或者说，它们的逼近误差能否与最佳逼近误差相比拟？
\end{enumerate}

本文的主线即是围绕这两个问题展开。Korovkin 定理给出了第一个问题的答案：对于线性正算子，只需验证三个简单的“测试函数”（$1, x, x^2$ 或 $1, \cos t, \sin t$）的逼近行为，就能保证对一切连续函数的一致收敛性。这一判别法极为简洁，它将众多经典算子的收敛性证明统一到了同一个框架之下。

第二个问题则需要将构造性逼近拉回到最佳逼近的参照系中来回答。Lebesgue 不等式提供了连接二者的桥梁：任何线性算子的逼近误差，可以被最佳逼近误差乘以该算子的 Lebesgue 常数所控制。而 Jackson 定理进一步给出了最佳逼近误差的定量估计——它由函数的光滑性（连续模）所决定。综合二者，即可对具体算子的逼近阶给出精确刻画。

本文的组织结构如下：第二章介绍最佳逼近的存在性、唯一性与切比雪夫交错定理，确立逼近的“最优基准”；第三章阐述 Korovkin 定理及其在 Bernstein、Fejér 等经典算子上的应用，展示构造性逼近的统一收敛性判据；第四章引入 Lebesgue 不等式，将线性算子的误差纳入最佳逼近的框架中进行分析；第五章利用 Jackson 定理与 Zygmund 类，对逼近阶进行定量刻画。最后给出总结与展望。